% ============================================================
%  13_Method.tex – YÖNTEM
% ============================================================
\section{YÖNTEM}
\label{sec:yontem}

Bu bölümde Myora platformunun sistem mimarisi, kullanılan teknolojiler, veritabanı tasarımı, Edge Function tabanlı yapay zekâ servisleri ve geliştirme süreci detaylı olarak açıklanmaktadır.

\subsection{Sistem Mimarisi Genel Bakış}

Myora, dört katmanlı bir Backend-as-a-Service (BaaS) mimarisi üzerine inşa edilmiştir. Klasik bir Express sunucusu yerine \textbf{Supabase} platformu (auth, PostgreSQL, Edge Functions ve depolama) kullanılarak operasyonel yük en aza indirilmiş, geliştirme hızı artırılmıştır~\cite{Supabase2024}. Mimari katmanları şunlardır: (i)~Capacitor 7 tabanlı mobil sarmalayıcı katmanı, (ii)~React 18 SPA istemci katmanı, (iii)~Supabase Auth + PostgreSQL + Deno Edge Function arka uç katmanı ve (iv)~OpenAI GPT-4o destekli yapay zekâ analiz katmanı (Şekil~\ref{fig:architecture}).

\begin{figure}[H]
\centering
\fbox{\parbox{0.85\textwidth}{
\centering
\textbf{Myora -- 4 Katmanlı Sistem Mimarisi}\\[8pt]
\begin{tabular}{|p{0.8\textwidth}|}
\hline
\textbf{Mobil Katman} (Capacitor 7.x) \\
Android (Java) \mbox{$\mid$} iOS (gelecek hedef) \\
Eklentiler: Camera, Haptics, Keyboard, LocalNotifications, Preferences \\
\hline
\textbf{İstemci Katmanı} (React 18.3 SPA, Vite + SWC, TypeScript) \\
Pages \mbox{$\rightarrow$} Lazy Components \mbox{$\rightarrow$} Hooks \mbox{$\rightarrow$} Contexts \\
TanStack Query \mbox{$\mid$} Tailwind \mbox{$\mid$} shadcn/ui (Radix) \\
$\updownarrow$ Supabase JS SDK (REST + Realtime) \\
\hline
\textbf{Arka Uç (Supabase BaaS)} \\
Supabase Auth (e-posta/parola, JWT, RLS politikaları) \\
PostgreSQL 15 + 33 tablo + pgvector eklentisi \\
3 Deno Edge Function (analyze-health-data, generate-chat-response, generate-health-profile) \\
\hline
\textbf{Yapay Zekâ Katmanı} \\
OpenAI GPT-4o (chat/vision) \mbox{$\mid$} text-embedding-3-small (1536 boyut) \\
RAG: pgvector + ivfflat indeksi + cosine benzerliği \\
\hline
\end{tabular}
}}
\caption{Myora platformunun katmanlı sistem mimarisi}
\label{fig:architecture}
\end{figure}

\subsection{Teknoloji Yığını}

\subsubsection{Ön Yüz (Frontend) Teknolojileri}

Ön yüz katmanı, modern web standartlarına uygun olarak \texttt{package.json} dosyasında tanımlı bağımlılıklar kullanılarak geliştirilmiştir (Tablo~\ref{tab:frontend_stack}).

\begin{table}[H]
\centering
\caption{Ön yüz teknoloji yığını}
\label{tab:frontend_stack}
\begin{tabular}{|l|l|l|}
\hline
\textbf{Teknoloji} & \textbf{Versiyon} & \textbf{Kullanım Amacı} \\
\hline
React              & 18.3 & Bileşen tabanlı UI çerçevesi~\cite{React2024} \\
Vite + SWC plugin  & 5.4 / 4.x & Hızlı HMR ile geliştirme aracı \\
TypeScript         & 5.8 & Strict mod tip güvenliği \\
Tailwind CSS       & 3.4 & Yardımcı sınıf tabanlı stillendirme \\
Radix UI / shadcn/ui & --   & Erişilebilir, başsız UI bileşenleri \\
TanStack React Query & 5.x & Sunucu durum yönetimi ve önbellek \\
React Router DOM    & 6.x & İstemci tarafı yönlendirme (BrowserRouter / HashRouter) \\
Recharts            & 2.x & Veri görselleştirme (lazy yüklenir) \\
Lucide React        & 0.462 & SVG ikon kütüphanesi \\
react-hook-form     & 7.x & Form yönetimi \\
Zod                 & 3.x & Şema doğrulama \\
@supabase/supabase-js & 2.x & Auth/REST/Realtime istemcisi \\
\hline
\end{tabular}
\end{table}

\subsubsection{Arka Uç (Backend) Teknolojileri}

\begin{table}[H]
\centering
\caption{Arka uç teknoloji yığını}
\label{tab:backend_stack}
\begin{tabular}{|l|l|l|}
\hline
\textbf{Teknoloji} & \textbf{Versiyon} & \textbf{Kullanım Amacı} \\
\hline
Supabase Platform   & --    & BaaS (auth, DB, edge fn, storage)~\cite{Supabase2024} \\
PostgreSQL          & 15+   & İlişkisel veritabanı~\cite{PostgreSQL2024} \\
pgvector eklentisi  & --    & Vektör benzerlik araması (1536 boyut) \\
Deno (Edge Runtime) & v1.x  & Sunucusuz fonksiyon çalışma zamanı \\
Supabase Auth       & --    & E-posta/parola, JWT, RLS \\
Row Level Security  & --    & Kullanıcı bazlı veri izolasyonu \\
OpenAI Node SDK     & 6.x   & GPT-4o ve embedding API erişimi \\
Drizzle Kit         & 0.31  & Sadece şema tipleri (\texttt{db:push}) \\
\hline
\end{tabular}
\end{table}

\subsubsection{Mobil Katman}

Capacitor 7 ile aynı kod tabanından native Android (\texttt{com.myora.app}) paketi üretilmektedir~\cite{Capacitor2024}. Web çıktısı (\texttt{dist/public}) \texttt{npx cap sync android} komutu ile Android Studio projesine kopyalanmaktadır. \texttt{capacitor.config.ts} dosyasındaki temel eklentiler: \texttt{@capacitor/camera} (yemek ve tıbbi fotoğraf), \texttt{@capacitor/haptics} (dokunmatik geri bildirim), \texttt{@capacitor/keyboard} (klavye davranışı), \texttt{@capacitor/local-notifications} (bildirim zamanlama) ve \texttt{@capacitor/preferences} (anahtar-değer depolama).

\subsection{Dizin Yapısı ve Kod Organizasyonu}

Proje, paylaşılan tip tanımları (\texttt{shared/schema.ts}) ile istemci ve Edge Function'lar arasında köprü kuran bir monorepo olarak organize edilmiştir.

\begin{table}[H]
\centering
\caption{Proje dizin yapısı (üst seviye)}
\label{tab:dir_structure}
\begin{tabular}{|l|l|}
\hline
\textbf{Dizin / Dosya} & \textbf{İçerik} \\
\hline
\texttt{client/src/components/}    & 41 üst-düzey React bileşeni \\
\texttt{client/src/components/ui/} & 53 shadcn/ui temel bileşeni \\
\texttt{client/src/components/charts/} & Recharts sarmalayıcıları (lazy) \\
\texttt{client/src/contexts/}      & 3 React Context (Auth, HealthData, Language) \\
\texttt{client/src/hooks/}         & 7 özel kanca (useHealthScore, useStreak, vb.) \\
\texttt{client/src/lib/api/}       & 8 modüler API katmanı (auth, health, chat, ...) \\
\texttt{client/src/pages/}         & AppLayout, Auth, Onboarding, Index, NotFound \\
\texttt{client/src/services/}      & Bildirim ve native widget servisleri \\
\texttt{supabase/migrations/}      & 14 SQL migrasyon dosyası (1219 satır) \\
\texttt{supabase/functions/}       & 3 Edge Function dizini \\
\texttt{shared/schema.ts}          & Drizzle tip şeması (yalnızca tip referansı) \\
\texttt{scripts/rag/}              & Bilimsel kaynak yutma (ingest) betikleri \\
\texttt{android/}                  & Capacitor Android projesi \\
\hline
\end{tabular}
\end{table}

\subsection{Veritabanı Tasarımı}

Veritabanı şeması, \texttt{supabase/migrations/} dizinindeki 14 ayrı SQL migrasyon dosyasında tanımlanmıştır ve toplamda 33 tablo ile bunların ilişkilerini, indekslerini ve Row Level Security (RLS) politikalarını içermektedir. Tüm tablolar UUID birincil anahtar (\texttt{gen\_random\_uuid()}) kullanmakta ve \texttt{auth.users} tablosuna \texttt{ON DELETE CASCADE} ilişkisi ile bağlanmaktadır.

Tablo~\ref{tab:db_modules} migrasyon modülleri ve kapsadıkları tabloları özetlemektedir.

\begin{table}[H]
\centering
\caption{Migrasyon modülleri ve içerdikleri tablolar}
\label{tab:db_modules}
\small
\begin{tabular}{|l|p{0.55\textwidth}|c|}
\hline
\textbf{Migrasyon} & \textbf{Tablolar} & \textbf{Sayı} \\
\hline
\texttt{profiles}              & profiles                                                           & 1 \\
\texttt{core\_health}          & daily\_health\_logs, food\_entries, activity\_entries, voice\_entries, health\_documents & 5 \\
\texttt{streak\_achievements}  & achievements, user\_achievements, user\_streaks                    & 3 \\
\texttt{goals\_community}      & user\_goals, community\_posts, community\_comments, community\_likes & 4 \\
\texttt{family\_supplements}   & family\_members, family\_medications, family\_medication\_logs, supplements, supplement\_recommendations, supplement\_orders & 6 \\
\texttt{health\_profiles}      & health\_profiles                                                   & 1 \\
\texttt{chat}                  & chat\_conversations, chat\_messages                                & 2 \\
\texttt{analysis\_assets}      & blood\_tests, blood\_test\_results, medical\_photos                & 3 \\
\texttt{wearables}             & wearable\_devices, user\_devices, wearable\_data                   & 3 \\
\texttt{fasting}               & fasting\_logs                                                      & 1 \\
\texttt{rag\_chunks}           & scientific\_sources, scientific\_source\_chunks, rag\_query\_logs  & 3 \\
\texttt{supplement\_tracking}  & supplement\_intake\_logs                                           & 1 \\
\hline
\textbf{Toplam}                &                                                                    & \textbf{33} \\
\hline
\end{tabular}
\end{table}

\subsubsection{Temel Tablolar ve Alanları}

\paragraph{Kullanıcı Profilleri (\texttt{profiles}):}
\texttt{auth.users} tablosuyla 1:1 ilişkili olup PRIMARY KEY \texttt{id uuid} alanı \texttt{auth.users(id)}'e referans vermektedir. Yeni bir kullanıcı kayıt olduğunda \texttt{handle\_new\_user\_profile} adlı tetikleyici (trigger) ile otomatik oluşturulur. Alanlar: \texttt{email}, \texttt{display\_name}, \texttt{date\_of\_birth}, \texttt{gender}, \texttt{height\_cm}, \texttt{weight\_kg numeric(5,2)}, \texttt{activity\_level} (varsayılan \texttt{moderate}), \texttt{language} (\texttt{tr}/\texttt{en}), \texttt{timezone} (\texttt{Europe/Istanbul}), \texttt{onboarding\_completed}.

\paragraph{Günlük Sağlık Kayıtları (\texttt{daily\_health\_logs}):}
\texttt{(user\_id, date)} bileşik benzersizlik kısıtlamasıyla günlük özet veri tutar: \texttt{calories\_consumed}, \texttt{calories\_burned}, \texttt{calories\_target} (varsayılan 2100), \texttt{protein\_grams}, \texttt{protein\_target} (180g), \texttt{carbs\_grams}, \texttt{fat\_grams}, \texttt{water\_ml}, \texttt{water\_target} (2400 mL), \texttt{steps}, \texttt{steps\_target} (10\,000), \texttt{sleep\_minutes}, \texttt{sleep\_target} (480 dk), \texttt{stress\_level}, \texttt{mood\_score}.

\paragraph{Besin Girişleri (\texttt{food\_entries}):}
\texttt{meal\_type}, \texttt{food\_name}, \texttt{calories}, makro besinler (\texttt{protein\_grams}, \texttt{carbs\_grams}, \texttt{fat\_grams}, \texttt{fiber\_grams}), \texttt{serving\_size}, \texttt{image\_url}, \texttt{barcode}, \texttt{ai\_confidence numeric(5,2)} ve \texttt{source} (\texttt{manual} / \texttt{camera} / \texttt{voice} / \texttt{barcode}).

\paragraph{Kan Tahlilleri (\texttt{blood\_tests} / \texttt{blood\_test\_results}):}
Üst tablo OCR metni, AI özeti ve genel durum bilgilerini; alt tablo her bir belirteci (marker adı, değer, birim, referans aralığı, durum: \texttt{normal}/\texttt{low}/\texttt{high}/\texttt{critical}, kategori, AI yorumu) tutar.

\paragraph{Sağlık Profili (\texttt{health\_profiles}):}
Kullanıcı başına 1:1 olup AI tarafından üretilen kapsamlı analizi tutar: \texttt{bmi}, \texttt{bmr}, \texttt{tdee}, \texttt{body\_fat\_estimate}, \texttt{health\_score} (0--100), \texttt{risk\_factors} (JSONB), \texttt{strengths} (text[]), \texttt{improvement\_areas} (text[]), \texttt{nutrition\_plan} (JSONB), \texttt{exercise\_plan} (JSONB), \texttt{supplement\_recommendations} (JSONB), \texttt{ai\_generated\_summary}.

\paragraph{RAG Altyapısı:}
\texttt{scientific\_sources} (başlık, yazarlar, dergi, DOI, kategori, doğrulanmış mı), \texttt{scientific\_source\_chunks} (parça metni, \texttt{embedding extensions.vector(1536)}, anahtar kelimeler, dil) ve \texttt{rag\_query\_logs} (her sorgu için getirilen parça kimlikleri, benzerlik skorları, model, gecikme) tablolarından oluşmaktadır. \texttt{idx\_chunks\_embedding} indeksi \texttt{ivfflat} ve \texttt{vector\_cosine\_ops} ile 50 küme üzerinde tanımlıdır. \texttt{public.match\_source\_chunks(query\_embedding, match\_threshold, match\_count, ...)} adlı PostgreSQL fonksiyonu Edge Function'lar tarafından RPC ile çağrılarak benzerlik araması yapar.

\paragraph{Sohbet (\texttt{chat\_conversations} / \texttt{chat\_messages}):}
Mesaj tablosunda \texttt{role} (\texttt{user}/\texttt{assistant}), \texttt{content}, \texttt{citations} (JSONB), \texttt{suggested\_follow\_ups} (text[]) ve \texttt{topic} alanları bulunmaktadır.

\paragraph{Row Level Security:}
Tüm hassas tablolar (\texttt{profiles}, \texttt{daily\_health\_logs}, \texttt{food\_entries}, \texttt{blood\_tests}, \texttt{health\_profiles}, vb.) için \texttt{enable row level security} ile RLS aktifleştirilmiş ve \texttt{(select auth.uid()) = user\_id} koşulu ile yalnızca veri sahibinin SELECT/INSERT/UPDATE/DELETE yapabileceği politikalar tanımlanmıştır. Bilimsel kaynaklar paylaşımlı bilgi tabanı olduğundan tüm kimlik doğrulanmış kullanıcılar tarafından okunabilir; yutma (ingestion) işlemleri \texttt{service\_role} ile yapılır.

\subsection{İstemci API Katmanı}

İstemci tarafı API katmanı \texttt{client/src/lib/api/} altında 8 ayrı modül halinde organize edilmiştir; toplam 2954 satır TypeScript koddan oluşmaktadır. Her modül namespaced bir nesne olarak dışa aktarılmaktadır (Tablo~\ref{tab:api_modules}).

\begin{table}[H]
\centering
\caption{İstemci API modülleri ve sorumlulukları}
\label{tab:api_modules}
\begin{tabular}{|l|c|p{0.45\textwidth}|}
\hline
\textbf{Modül} & \textbf{Satır} & \textbf{Sorumluluk} \\
\hline
\texttt{auth.ts}        & 97   & Kayıt, giriş, çıkış, profil güncelleme \\
\texttt{health.ts}      & 978  & Günlük log, besin, aktivite, kan, belge, fotoğraf, ses, oruç, hedef, başarım \\
\texttt{chat.ts}        & 258  & Konuşma listesi, mesaj gönderme, RAG yanıtı, sağlık profili oluşturma \\
\texttt{supplement.ts}  & 224  & Katalog, AI önerileri, sipariş ve alım kayıtları \\
\texttt{community.ts}   & 158  & Gönderi/yorum/beğeni \\
\texttt{family.ts}      & 251  & Aile üyeleri, ilaç ve kayıtları, uyum oranları \\
\texttt{wearable.ts}    & 129  & Cihaz katalogu, eşleme, veri senkronizasyonu \\
\texttt{notification.ts}& 52   & Tercih yönetimi ve haftalık görevler \\
\texttt{\_common.ts}    & 117  & Supabase istemci, oturum yardımcıları, hata sarmalayıcıları \\
\texttt{types.ts}       & 657  & Paylaşılan tip tanımları \\
\hline
\textbf{Toplam}         & \textbf{2921} & \\
\hline
\end{tabular}
\end{table}

\subsubsection{Kimlik Doğrulama ve Oturum Yönetimi}

Kimlik doğrulama, Supabase Auth tarafından sağlanmaktadır. Akış şu şekildedir:

\begin{enumerate}
    \item Kullanıcı \texttt{authAPI.signUp(email, password, displayName)} çağrısı ile kayıt olur. Bu çağrı Supabase \texttt{auth.signUp} metodunu kullanır.
    \item \texttt{auth.users} tablosuna kayıt yapıldığında \texttt{handle\_new\_user\_profile} tetikleyicisi \texttt{public.profiles} tablosuna ilgili satırı oluşturur.
    \item Giriş için \texttt{authAPI.signIn} kullanılır; başarılı oturumla birlikte JWT erişim simgesi (access token) üretilir ve istemci tarafında \texttt{onAuthStateChange} dinleyicisi ile takip edilir.
    \item Tüm sonraki API çağrıları, Supabase JS SDK üzerinden \texttt{Authorization: Bearer <jwt>} başlığı ile yapılır; veriye erişim RLS politikaları ile sınırlandırılır.
    \item \texttt{ProtectedRoute} bileşeni \texttt{AuthContext}'i tüketerek korumalı rotaları yetkisiz kullanıcılardan yönlendirir.
\end{enumerate}

\subsection{Yapay Zekâ Servisleri (Supabase Edge Functions)}

Myora'nın yapay zekâ katmanı, Deno çalışma zamanında çalışan üç ayrı Supabase Edge Function ile gerçekleştirilmiştir (Tablo~\ref{tab:ai_services}). Tüm fonksiyonlar OpenAI GPT-4o modelini kullanmakta, yapılandırılmış JSON çıktı formatı (\texttt{response\_format: \{type: "json\_object"\}}) ile yanıt vermekte ve ALLOWED\_ORIGINS yapılandırması ile CORS güvenliği sağlamaktadır.

\begin{table}[H]
\centering
\caption{Edge Function tabanlı yapay zekâ servisleri}
\label{tab:ai_services}
\begin{tabular}{|l|c|p{0.45\textwidth}|}
\hline
\textbf{Edge Function} & \textbf{Satır} & \textbf{Görev} \\
\hline
\texttt{analyze-health-data}     & 203 & 6 alt-tip için tek uç nokta: \texttt{food-image}, \texttt{blood-test}, \texttt{document}, \texttt{medical-photo}, \texttt{voice-entry}, \texttt{voice-food} \\
\texttt{generate-chat-response}  & 285 + 843 (lib) & 8 aşamalı RAG pipeline (Auth $\to$ Safety $\to$ Topic $\to$ Context $\to$ Retrieval $\to$ Prompt $\to$ Generation $\to$ Logging) \\
\texttt{generate-health-profile} & 259 & BMI, BMR (Mifflin--St Jeor), TDEE, sağlık skoru ve plan üretimi \\
\hline
\end{tabular}
\end{table}

\subsubsection{analyze-health-data: Çok Kipli Analiz}

\texttt{type} alanına göre dinamik bir prompt seçilerek GPT-4o'ya çağrı yapılır; OpenAI anahtarı yoksa veya çağrı başarısızsa servis, alanın varsayılan/yedek (fallback) bir JSON nesnesi döndürür. Örnek alt-tipler:
\begin{itemize}
    \item \texttt{food-image} -- base64 görüntü ile besin adı, kalori, porsiyon, makrolar (protein/carbs/fat/fiber, gram + yüzde) ve sağlık notları üretir.
    \item \texttt{medical-photo} -- 5 fotoğraf tipi (\texttt{urine}, \texttt{stool}, \texttt{tongue}, \texttt{eye}, \texttt{skin}) için renk, gözlem, endişe ve aciliyet (\texttt{low}/\texttt{medium}/\texttt{high}/\texttt{urgent}) üretir.
    \item \texttt{blood-test} -- OCR metninden belirteçleri çıkararak kategori, durum ve öneri listesi sunar.
    \item \texttt{voice-entry} / \texttt{voice-food} -- transkripsiyondan duygu, ruh hali, anahtar kelimeler veya yenen yiyecekleri çıkarır.
\end{itemize}

\subsubsection{generate-chat-response: 8 Aşamalı RAG Pipeline}

RAG chatbot, Şekil~\ref{fig:rag}'te gösterilen sekiz aşamalı bir boru hattıyla çalışmaktadır~\cite{Lewis2020RAG}. Aşamalar \texttt{supabase/functions/generate-chat-response/lib/} altındaki beş modül (\texttt{safety.ts}, \texttt{topic.ts}, \texttt{context.ts}, \texttt{retrieval.ts}, \texttt{prompt.ts}) tarafından gerçekleştirilmektedir.

\begin{figure}[H]
\centering
\fbox{\parbox{0.85\textwidth}{
\centering
\textbf{RAG Chatbot Veri Akışı (8 Aşama)}\\[8pt]
1. Kullanıcı sorgusu + auth (JWT doğrulama) \\
$\downarrow$ \\
2. \textbf{Safety} -- 7 yüksek riskli intent örüntüsü (göğüs ağrısı, intihar, anafilaksi, vb.) \\
$\downarrow$ \\
3. \textbf{Topic} -- konu sınıflandırması (\texttt{nutrition}, \texttt{sleep}, \texttt{exercise}, ...) \\
$\downarrow$ \\
4. \textbf{Context} -- profil, kan, oruç, su, son besinler (paralel sorgular) \\
\quad $\parallel$ \quad \\
5. \textbf{Retrieval} -- 1536 boyutlu embedding + \texttt{match\_source\_chunks} RPC (\texttt{matchThreshold=0.3}, \texttt{matchCount=8}) \\
$\downarrow$ \\
6. \textbf{Prompt} -- rol + kurallar + bilimsel kaynaklar + kullanıcı bağlamı \\
$\downarrow$ \\
7. \textbf{Generation} -- GPT-4o JSON yanıtı: \{content, citations[], suggestedFollowUps[]\} \\
$\downarrow$ \\
8. \textbf{Logging} -- \texttt{rag\_query\_logs}'a sorgu, parça kimlikleri, benzerlik skorları ve gecikme yazılır \\
}}
\caption{Myora RAG chatbot 8-aşamalı boru hattı}
\label{fig:rag}
\end{figure}

\paragraph{Güvenlik Katmanı (Safety):}
\texttt{lib/safety.ts}, kardiyak acil, solunum acili, bilinç kaybı, intihar düşüncesi, hemoraji, anafilaksi ve zehirlenme olmak üzere yedi yüksek riskli desen içerir. Bu örüntüler tetiklendiğinde model, Türkiye için 112 (acil) ve 182 (intihar önleme) hatlarına yönlendirme yapan acil mesaj döndürür.

\paragraph{Geri Çağırma (Retrieval):}
Kullanıcı sorgusu önce \texttt{text-embedding-3-small} modeli ile 1536 boyutlu vektöre dönüştürülür, ardından PostgreSQL fonksiyonu \texttt{public.match\_source\_chunks} cosine benzerliği ile en alakalı 8 parçayı getirir. Yalnızca \texttt{is\_verified = true} olan kaynaklar değerlendirilir.

\subsubsection{generate-health-profile: Skor ve Plan Üretimi}

\texttt{generate-health-profile} servisi BMI, Mifflin--St Jeor formülüyle BMR, aktivite çarpanlarıyla TDEE hesaplar. Aktivite çarpanları: \texttt{sedentary} 1.2, \texttt{light} 1.375, \texttt{moderate} 1.55, \texttt{active} 1.725, \texttt{very\_active} 1.9. 0--100 arası sağlık skoru, BMI normal aralık (+10), düzenli öğün takibi (+10), su hedefine ulaşma (+10), aktivite (+10), uyku (+5) gibi faktörlerden oluşur. AI çağrısı başarısız olduğunda da kural-tabanlı yedek profil (fallback) döndürülmesi sağlanmıştır.

\subsection{Ön Yüz Bileşen Mimarisi}

\texttt{client/src/components/} dizininde 41 üst-düzey React bileşeni, \texttt{client/src/components/ui/} altında 53 shadcn/ui temel bileşeni yer almaktadır. Tüm sayfa düzeyinde bileşenler \texttt{App.tsx} içinde \texttt{React.lazy()} ile tembel yüklenmektedir; bu sayede ilk yükleme paket boyutu küçük tutulmuştur.

\subsubsection{Sayfa Düzeni}

\texttt{pages/AppLayout.tsx} korumalı bir kabuktur; alt rotalar (\texttt{HealthDashboard}, \texttt{AnalysisHub}, \texttt{Boost}, \texttt{ExpertChat} vb.) iç içe rotalar olarak yerleştirilmiştir. Web ortamında \texttt{BrowserRouter}, native ortamda (\texttt{Capacitor.isNativePlatform()}) \texttt{HashRouter} kullanılır. \texttt{ErrorBoundary} sınıf bileşeni rota seviyesinde tüm beklenmeyen hataları yakalayarak \enquote{Tekrar dene} aksiyonlu kullanıcı dostu bir UI sunar.

\subsubsection{Durum Yönetimi}

Uygulama durumu üç bağlam sağlayıcısı ve TanStack Query ile yönetilmektedir:

\begin{itemize}
    \item \textbf{AuthContext:} Supabase oturumunu (\texttt{user}, \texttt{session}) ve \texttt{signIn}/\texttt{signUp}/\texttt{signOut}/\texttt{updateProfile} işlemlerini yönetir.
    \item \textbf{HealthDataContext:} İstemci tarafı önbellek (yemek analizleri, sesli girişler, aktiviteler, belgeler, tıbbi fotoğraflar, günlük hedefler, su tüketimi). Günlük kalori ve protein toplamlarını \texttt{useMemo} ile otomatik hesaplar.
    \item \textbf{LanguageContext:} \texttt{tr}/\texttt{en} tercihi, \texttt{t} fonksiyonu ile çeviri sözlüğü. Marka isimleri (\enquote{Myora}, \enquote{My Score}) iki dilde de İngilizce kalır.
\end{itemize}

Sunucu durumu ise \texttt{TanStack React Query} ile yönetilmekte; tüm sorgu anahtarları \texttt{client/src/lib/queryKeys.ts} dosyasında merkezileştirilmiştir. \texttt{staleTime.static} ve \texttt{staleTime.dynamic} sabitleri ile farklı veri türlerine farklı önbellek ömrü uygulanmıştır.

\subsubsection{Kanca (Hook) Kütüphanesi}

\texttt{client/src/hooks/} altında uygulamanın yedi özel kancası bulunmaktadır:
\begin{itemize}
    \item \texttt{useHealthScore}, \texttt{useStreak} -- TanStack Query ile sağlık skoru ve seri takibi.
    \item \texttt{useHealthData} -- HealthDataContext sarmalayıcısı.
    \item \texttt{use-mobile}, \texttt{useKeyboard} -- mobil platform algılama ve klavye davranışı.
    \item \texttt{useTabNavigate} -- sekme bazlı navigasyon.
    \item \texttt{use-toast} -- shadcn toast bildirimleri.
\end{itemize}

\subsection{Form Doğrulama ve i18n}

Tüm formlar \texttt{react-hook-form} ile yönetilmekte, \texttt{@hookform/resolvers/zod} adaptörü ile Zod şemalarına bağlanmaktadır. Şemalar \texttt{client/src/lib/validations/} altında \texttt{auth.ts} ve \texttt{forms.ts} dosyalarında tanımlıdır. Çoklu dil desteği için \texttt{LanguageContext} 200'den fazla çeviri anahtarı içermekte, tip-güvenli olarak \texttt{en} ve \texttt{tr} sözlüklerinde aynı anahtar kümesi zorlanmaktadır.

\subsection{Bilimsel Kaynak Yutma (RAG Ingest) Pipeline}

\texttt{scripts/rag/} dizini iki TypeScript betiği içermektedir:
\begin{itemize}
    \item \texttt{seed-sources.ts} (581 satır) -- önceden tanımlı bilimsel kaynak meta verilerini \texttt{scientific\_sources} tablosuna yükler.
    \item \texttt{ingest-urls.ts} (543 satır) -- URL'lerden makale içeriğini çekip parçalara böler, \texttt{text-embedding-3-small} ile gömme vektörleri üretir ve \texttt{scientific\_source\_chunks} tablosuna yazar.
\end{itemize}
Bu betikler \texttt{npm run rag:seed} ve \texttt{npm run rag:ingest} komutları ile çalıştırılmaktadır.

\subsection{Geliştirme ve Dağıtım Süreci}

Proje, aşamalı bir geliştirme sürecine tabi tutulmuştur:

\begin{enumerate}
    \item \textbf{Ocak 2025:} İlk prototip Lovable platformunda Supabase tabanlı olarak geliştirilmiştir.
    \item \textbf{Ocak 2026:} Replit ortamına taşınmış ve geçici olarak Express tabanlı bir vites/middleware katmanı eklenmiştir; sonrasında bu katman kaldırılarak doğrudan Supabase Edge Function'larına geçiş yapılmıştır.
    \item \textbf{Şubat 2026:} Capacitor 7 entegrasyonu ile Android paketleme tamamlanmıştır.
    \item \textbf{Şubat 2026:} 5 sekmeli alt navigasyon, lazy yükleme ve ErrorBoundary eklenerek bilgi mimarisi ve performans iyileştirmeleri yapılmıştır.
    \item \textbf{Mart 2026:} \texttt{rag\_chunks} ve \texttt{supplement\_tracking} migrasyonları, \texttt{ivfflat} indeksi ve RLS audit düzeltmeleri uygulanmıştır.
\end{enumerate}

\paragraph{Ortam Değişkenleri:}
\texttt{VITE\_SUPABASE\_URL}, \texttt{VITE\_SUPABASE\_PUBLISHABLE\_KEY} (istemci); \texttt{OPENAI\_API\_KEY}, \texttt{SUPABASE\_URL}, \texttt{SUPABASE\_SERVICE\_ROLE\_KEY}, \texttt{ALLOWED\_ORIGINS} (Edge Function gizli değerleri).

\paragraph{Dağıtım Komutları:}
\texttt{npm run supabase:db:push} ile şema, \texttt{npm run supabase:functions:deploy} ile üç Edge Function tek seferde dağıtılır. Web istemci için Netlify (\texttt{netlify.toml}), Vercel (\texttt{vercel.json}) ve Railway (\texttt{railway.json}) yapılandırmaları hazırdır.

\subsection{Kod Kalitesi ve Tip Güvenliği}

TypeScript \texttt{strict} modu açıktır. ESLint kuralları \texttt{eslint.config.js} dosyasında flat config formatında tanımlanmıştır. \texttt{npx tsc --noEmit} ile tip kontrolü, \texttt{npm run lint} ile statik analiz çalıştırılmaktadır. Tüm bileşenlerde \texttt{any}, \texttt{getattr}/\texttt{setattr} kalıpları yerine açık tip tanımları tercih edilmiştir.
